\documentclass[aspectratio=169]{beamer}

\usetheme{default}
\usecolortheme{default}
\setbeamertemplate{navigation symbols}{}
\setbeamertemplate{footline}[frame number]

\usepackage[utf8]{inputenc}
\usepackage[T1]{fontenc}
\usepackage[spanish]{babel}
\usepackage{amsmath}
\usepackage{ragged2e}
\usepackage{xcolor}
\usepackage{graphicx} % For vertical centering

% --------------------------
% Colors & Style (unchanged)
% --------------------------
\definecolor{YPFBlue}{RGB}{0,51,102}
\definecolor{LightBlue}{RGB}{102,178,255}
\definecolor{SoftGray}{RGB}{90,90,90}
\definecolor{VeryLightGray}{RGB}{240,242,245}

\setbeamercolor{title}{fg=YPFBlue}
\setbeamercolor{frametitle}{fg=YPFBlue}
\setbeamercolor{normal text}{fg=black,bg=white}
\setbeamercolor{block title}{bg=VeryLightGray,fg=YPFBlue}
\setbeamercolor{block body}{bg=white}


\setbeamerfont{title}{series=\bfseries,size=\LARGE}
\setbeamerfont{frametitle}{series=\bfseries,size=\Large}

\newcommand{\highlight}[1]{\textcolor{YPFBlue}{\textbf{#1}}}
\newcommand{\soft}[1]{\textcolor{SoftGray}{#1}}

% --------------------------
% Title info
% --------------------------
\title{Datos, econometría e inferencia causal para decisiones de negocio}
\author{
Santiago Cerutti \and Guillermo Falcone \and Francisco Pizzi \and Jorge Puig
}
\date{}

\begin{document}

% ==================================================
% Slide 1: Team / Cover (Single Column Layout)
% ==================================================
\begin{frame}[plain]
    \centering
    \vspace{1cm}
    
    {\color{YPFBlue}\bfseries\LARGE Datos, econometría e inferencia causal\\[0.2cm]
    para decisiones de negocio}

    \vspace{0.7cm}

    {\large \soft{Una introducción a nuestro método de trabajo}}

    \vfill

    \begin{block}{\large Equipo}
        \large
        \quad\textbullet\quad Santiago Cerutti \\ \quad\textbullet\quad Guillermo Falcone \\ \quad\textbullet\quad Francisco Pizzi \\ \quad\textbullet\quad Jorge Puig
    \end{block}
    
    \vfill
\end{frame}

% ==================================================
% Slide 2: The Challenge of Deciding in the Real World (Revised, more sober)
% ==================================================
\begin{frame}{Cómo medir el impacto real de las decisiones de un negocio?}

\vspace{0.5cm}

\begin{center}
\Large
Las decisiones de una empresa no se toman en un laboratorio, sino en un \soft{entorno complejo y dinámico donde múltiples factores influyen simultaneamente sobre el resultado final.}
\end{center}

\vspace{1cm}

\begin{block}{El problema central: separar la señal del ruido}
\justifying
Cuando se toma una decisión, los resultados que vemos son una mezcla del impacto real de dicha acción y del "ruido" de fondo (la competencia, la economía, el clima, otras decisiones simultáneas, etc.). Una simple correlación puede ser muy engañosa.
\vspace{0.3cm}
\begin{center}
\Large\highlight{Entonces, ¿cómo aislamos el verdadero efecto de una decisión?}
\end{center}
\end{block}

\vfill

\end{frame}

% ==================================================
% Slide 4: Our Answer: The Counterfactual (Old Slide 3, renumbered & retitled)
% ==================================================
\begin{frame}{Nuestra respuesta: la ciencia de la inferencia causal}
\begin{center}
{\Large \soft{Para aislar el efecto, debemos comparar lo que pasó con lo que hubiese pasado.}}
\end{center}

\vspace{0.5cm}

\begin{center}
\Large
\[
\text{Resultado observado}
\quad \highlight{\text{vs.}} \quad
\underbrace{
    \begin{tabular}{c}
    Resultado en un mundo \\
    donde la acción no ocurrió
    \end{tabular}
}_{\text{\normalsize\highlight{El contrafactual}}}
\]
\end{center}

\vspace{0.5cm}

\begin{block}{Nuestro trabajo: construir ese mundo}
\justifying
Usando los propios datos de la empresa, construimos una \highlight{representación científica} de ese escenario contrafactual. Esto nos permite obtener una estimación del \highlight{efecto causal} real de la decisión.
\end{block}

\end{frame}

% ==================================================
% Slide 4: The Example - Part 1: The Question (NEW, split slide)
% ==================================================
\begin{frame}{Un ejemplo: ¿funcionó una promoción?}

\begin{block}{La pregunta de negocio}
\justifying
YPF lanza una promoción en 100 de sus estaciones. El objetivo es claro: queremos saber si funcionó y, sobre todo, \highlight{cuánto valor generó}.
\vspace{0.2cm}
\begin{center}
    \large
    \highlight{¿Cuál fue el impacto real de la promoción en las ventas?}
\end{center}
\end{block}

\begin{block}{El análisis "naive" (y por qué falla)}
\justifying
La primera tentación es comparar las ventas de las 100 estaciones \textit{con promoción} contra las ventas del resto.
\vspace{0.3cm}

\textit{\soft{Pero este enfoque es engañoso. Las estaciones no se eligen al azar; quizás la promoción se aplicó en las que ya eran más grandes o estaban en zonas de mayor poder adquisitivo. Comparar grupos que no son comparables nos lleva a conclusiones erróneas.}}
\end{block}

\end{frame}

% ==================================================
% Slide 5: The Example - Part 2: The Solution (NEW, split slide)
% ==================================================
\begin{frame}{Un ejemplo: nuestra solución}

\begin{block}{Paso 1: Construir un grupo de control a medida}
\justifying
En lugar de comparar contra "todas las demás", usamos los datos para encontrar \highlight{"estaciones gemelas"}. Para cada estación con promoción, identificamos un grupo de estaciones muy parecidas (en ventas previas, ubicación, tamaño, etc.) pero \highlight{que no tuvieron la promoción}.
\end{block}

\begin{block}{Paso 2: Comparar la evolución de ambos grupos}
\justifying
Ahora sí tenemos una comparación justa. Medimos el cambio en ventas \highlight{antes y después} de la promoción en ambos grupos y comparamos esas trayectorias.
\begin{center}
\large
\[
\underbrace{\Delta \text{Ventas}}_{\text{con promoción}}
\quad \text{vs.} \quad
\underbrace{\Delta \text{Ventas}}_{\text{\highlight{grupo de control}}}
\]
\end{center}
\justifying
La diferencia entre ambos es nuestra estimación del \highlight{efecto causal real}.
\end{block}

\end{frame}

% ==================================================
% Slide 6: How We Do It (Single Column Layout)
% ==================================================
\begin{frame}{¿Cómo lo hacemos? Un enfoque científico y a medida}

\begin{block}{Principio de trabajo}
\justifying
Nos enfocamos en \highlight{una pregunta específica a la vez}. Cada pregunta exige un diseño de análisis propio, porque los datos y el contexto son únicos. No usamos "modelos enlatados".
\end{block}

\begin{block}{Nuestra caja de herramientas}
\justifying
Aplicamos técnicas de la econometría de evaluación de impacto. El nombre es académico, pero la idea es simple: usar la herramienta correcta para cada trabajo. Algunas de ellas son:
\begin{itemize}
    \item Diseños de Diferencias en Diferencias
    \item Control Sintético (y derivados)
    \item Variables Instrumentales y Diseños de Regresión por Discontinuidad
    \item \highlight{Diseños ad-hoc que se ajustan a la estructura de los datos de YPF}
\end{itemize}
\end{block}

\vfill
\begin{center}
\Large
\soft{El foco es la} \highlight{identificación causal}, \soft{no la simple correlación o el pronóstico.}
\end{center}
\end{frame}

% ==================================================
% Slide 7: Why this matters (Single Column Layout)
% ==================================================
\begin{frame}{¿Por qué es valioso este enfoque?}

\begin{block}{Permite tomar mejores decisiones}
\justifying
Buscamos construir una base de conocimiento para la estrategia de la empresa. Es un activo que permite aprender del pasado para optimizar el futuro. Los resultados concretos son:
\begin{itemize}
    \item \highlight{Cuantificar} el ROI real de las acciones de negocio.
    \item \highlight{Reducir} la incertidumbre en la planificación y la asignación de recursos.
    \item \highlight{Separar} la intuición de la evidencia empírica para entender qué funciona y qué no.
\end{itemize}
\end{block}

\begin{block}{Cómo se integra en la empresa}
\justifying
Este conocimiento \highlight{no reemplaza} la experiencia ni la intuición del equipo. Las \highlight{complementa} y las potencia con evidencia medible, rigurosa y auditable.
\end{block}
\end{frame}

% ==================================================
% Slide 8: Closing (Single Column Layout)
% ==================================================
\begin{frame}{Nuestro compromiso: rigurosidad y pragmatismo}
\begin{center}
\Large
\soft{Somos un equipo de economistas aplicados.}
\vspace{0.4cm}

{\Large \highlight{Unimos la rigurosidad académica con el pragmatismo del negocio.}}
\end{center}

\vspace{1cm}

\begin{block}{Qué entregamos}
\justifying
El producto no es un reporte genérico, sino una \highlight{serie de respuestas precisas a preguntas concretas}. La acumulación de estas respuestas genera un entendimiento estratégico profundo sobre las palancas de valor de la empresa.
\end{block}

\vfill

\begin{center}
\LARGE
\soft{Transformamos sus datos en} \highlight{evidencia para decidir.}
\end{center}
\end{frame}

\end{document}